\section{Matrices}
\label{sec:matrices}

\noindent\textbf{Subject:} Linear Algebra for ML (3Blue1Brown) \\
\textbf{Status:} growing

\subsection*{Summary}
A matrix encodes a linear transformation. Matrix multiplication represents the composition of two transformations applied in sequence (read right to left).

\subsection{Linear Transformation}

\begin{defbox}
Transformation is a fancy word for \textbf{function} --- it takes an input and spits out an output.
\end{defbox}

\paragraph{What makes it \emph{Linear}?} Two properties:
\begin{enumerate}[nosep]
  \item All lines must remain lines \textbf{without getting curved}
  \item The \textbf{origin must remain fixed} in place
\end{enumerate}

$\rightarrow$ So a Linear Transformation is one that \textbf{keeps grid lines parallel and evenly spaced}.

\paragraph{How to describe one numerically?}
\begin{itemize}[nosep]
  \item We need to record where the basis vectors $\hat{i}$ and $\hat{j}$ land
  \item To find where \textbf{any} vector lands, use:
\end{itemize}

\[
\begin{bmatrix} x \\ y \end{bmatrix}
\rightarrow
x\begin{bmatrix} 1 \\ 0 \end{bmatrix} + y\begin{bmatrix} 0 \\ 1 \end{bmatrix}
\Rightarrow
x\begin{bmatrix} 1 \\ -2 \end{bmatrix} + y\begin{bmatrix} 3 \\ 0 \end{bmatrix}
= \begin{bmatrix} 1x+3y \\ -2x+0y \end{bmatrix}
\]

You can use a \textbf{$2\times2$ matrix} for $\hat{i}$, $\hat{j}$:

\[
\begin{bmatrix} a & 3 \\ -2 & 0 \end{bmatrix}
\]

To find where a transformation takes the vector $\begin{bmatrix} 5 \\ 7 \end{bmatrix}$, apply:

\[
5\begin{bmatrix} 1 \\ -2 \end{bmatrix} + 7\begin{bmatrix} 3 \\ 0 \end{bmatrix}
\]

\paragraph{What is ``Shear''?}
\begin{defbox}
It's a kind of transformation to the basis vectors (skewing the grid).
\end{defbox}

\paragraph{What is ``Composition''?}
\begin{defbox}
Rotation and a shear combined.
\end{defbox}

To know where the vector ends up:

\[
\underbrace{\begin{bmatrix} 1 & 1 \\ 0 & 1 \end{bmatrix}}_{\text{Shear}}
\left(
\underbrace{\begin{bmatrix} 0 & -1 \\ 1 & 0 \end{bmatrix}}_{\text{Rotation}}
\begin{bmatrix} x \\ y \end{bmatrix}
\right)
=
\underbrace{\begin{bmatrix} 1 & -1 \\ 1 & 0 \end{bmatrix}}_{\text{Composition}}
\begin{bmatrix} x \\ y \end{bmatrix}
\]

\begin{notebox}
Read \textbf{right to left}.
\end{notebox}

\subsection{Matrix Multiplication}

To multiply $M_1$ and $M_2$:

\[
\begin{bmatrix} a & b \\ c & d \end{bmatrix}
\begin{bmatrix} e & f \\ g & h \end{bmatrix}
=
\begin{bmatrix} ae+bg & af+bh \\ ce+dg & ac+dh \end{bmatrix}
\]

\begin{notebox}
\begin{itemize}[nosep]
  \item $M_1 \cdot M_2 \neq M_2 \cdot M_1$ (not commutative)
  \item But: $(AB)C = A(BC)$ (associative \checkmark)
\end{itemize}
\end{notebox}

\subsection*{Related Topics}
\begin{itemize}[nosep]
  \item Section~\ref{sec:vectors} --- Vectors
  \item Section~\ref{sec:determinant} --- Determinant
  \item Section~\ref{sec:linear-systems} --- Linear System of Equations
\end{itemize}

\bigskip
\noindent\textit{Tags: linear-algebra, matrices, transformation, math, phase-1}
